\documentclass[12pt,a4paper]{article}

% ── Packages ──────────────────────────────────────────────────────────────
\usepackage[utf8]{inputenc}
\usepackage[T1]{fontenc}

\usepackage[margin=2.5cm]{geometry}
\usepackage{setspace}
\usepackage{graphicx}
\usepackage{booktabs}
\usepackage{tabularx}
\usepackage{array}
\usepackage{amsmath,amssymb}
\usepackage[protrusion=false,expansion=false]{microtype}
\usepackage{hyperref}
\usepackage{xcolor}
\usepackage{natbib}
\usepackage{caption}
\usepackage{float}
\usepackage{lineno}
\usepackage{abstract}
\usepackage{titlesec}
\usepackage{enumitem}

% ── Hyperref setup ────────────────────────────────────────────────────────
\hypersetup{
  colorlinks=true,
  linkcolor=blue!70!black,
  citecolor=blue!70!black,
  urlcolor=blue!70!black,
  pdftitle={Beyond the Flight Recorder: A Cryptographic Biometric Framework for
            Bridging Individual and Organizational Aviation Risk},
  pdfauthor={Mete CANTEKİN},
  pdfkeywords={Aviation Safety, CDSA, Cryptographic Privacy, Shamir Secret Sharing,
               Differential Privacy, Wearable Biometric Sensors}
}

% ── Typography ────────────────────────────────────────────────────────────
\onehalfspacing
\setlength{\parindent}{1.25cm}
\setlength{\parskip}{0pt}

% ── Section formatting ────────────────────────────────────────────────────
\titleformat{\section}{\normalfont\large\bfseries}{\thesection.}{0.5em}{}
\titleformat{\subsection}{\normalfont\normalsize\bfseries}{\thesubsection.}{0.5em}{}
\titleformat{\subsubsection}{\normalfont\normalsize\bfseries}{\thesubsubsection.}{0.5em}{}
\titlespacing{\section}{0pt}{18pt}{6pt}
\titlespacing{\subsection}{0pt}{12pt}{4pt}
\titlespacing{\subsubsection}{0pt}{10pt}{2pt}

% ── Caption setup ─────────────────────────────────────────────────────────
\captionsetup{font=small, labelfont=bf, justification=justified, singlelinecheck=false}

% ── arXiv preprint note ───────────────────────────────────────────────────
\newcommand{\preprintnote}{%
  \begin{center}
    \fbox{\parbox{0.85\textwidth}{\centering\small
      \textbf{Preprint} --- Submitted to \textit{Safety Science} (Elsevier, Q1).\\
      This version has not yet undergone peer review.\\
      \href{https://orcid.org/0009-0001-6990-6340}{ORCID: 0009-0001-6990-6340}
      $\cdot$ Istanbul, 2026
    }}
  \end{center}
  \vspace{1em}
}

% ══════════════════════════════════════════════════════════════════════════
\begin{document}
% ══════════════════════════════════════════════════════════════════════════

% ── Title ─────────────────────────────────────────────────────────────────
\begin{titlepage}
\begin{center}
\vspace*{2cm}

{\LARGE\bfseries Beyond the Flight Recorder:\\[6pt]
A Cryptographic Biometric Framework for Bridging\\[6pt]
Individual and Organizational Aviation Risk}

\vspace{1.5cm}

{\large Mete CANTEKİN}\\[4pt]
{\normalsize Istanbul Beykent University, Department of Computer Engineering (PhD)}\\[2pt]
{\normalsize ORCID: \href{https://orcid.org/0009-0001-6990-6340}{0009-0001-6990-6340}}\\[2pt]
{\normalsize Supervisor: Assoc.\ Prof.\ Dr.\ İbrahim Furkan İNCE}

\vspace{1.5cm}

\rule{0.7\textwidth}{0.4pt}

\vspace{0.8cm}

{\small\itshape Preprint submitted to \textbf{Safety Science} (Q1, Elsevier)\\
Not yet peer-reviewed --- Istanbul, 2026}

\vspace{0.8cm}
\rule{0.7\textwidth}{0.4pt}

\end{center}
\end{titlepage}

% ── Abstract ──────────────────────────────────────────────────────────────
\begin{abstract}
Every pilot error is an organizational failure \citep{reason2000,dekker2007}.
Aviation safety has achieved remarkable maturity through the systematic development
of reactive and proactive paradigms over the past half century. Yet the mechanical
telemetry provided by flight data recorders (FDR) and cockpit voice recorders (CVR)
cannot capture the cognitive and physiological state of the pilot at the moment of
decision-making, creating a structural diagnostic gap in accident investigations.
This study presents a two-phase cryptographic data processing architecture developed
within the Complementary Diagnostic Safety Approach (CDSA) framework to address this gap.
In the first phase, HRV, SpO\textsubscript{2}, and GSR data from smart pilot watches
are encrypted at the device level using AES-256-GCM; the master key is then split
into three shares via Shamir Secret Sharing and distributed to independent institutional parties.
The 2-of-3 threshold scheme ensures that individual biometric data can only be accessed
following an accident report with simultaneous authorization from at least two parties.
In the second phase, macro biometric data anonymized through differential privacy
mechanisms ($\varepsilon=1.0$, $k$-anonymity $k=5$) are integrated into Safety
Management Systems (SMS) to enable proactive risk analysis.
Validated at prototype level in Python, the architecture is compliant with the
privacy-by-design principle under GDPR Article 25 and with Just Culture philosophy.
All reported findings are derived from a literature-calibrated synthetic dataset
($n=1{,}179$); real operational validation constitutes the primary agenda of future research.
The study also proposes a Biometric Flight Data Declaration within ICAO Annex~13.
\end{abstract}

\noindent\textbf{Keywords:} Aviation Safety; Complementary Diagnostic Safety Approach (CDSA);
Cryptographic Privacy; Shamir Secret Sharing; Differential Privacy; Wearable Biometric Sensors

\vspace{1em}
\hrule
\vspace{1em}

% ══════════════════════════════════════════════════════════════════════════
\section{Introduction}
% ══════════════════════════════════════════════════════════════════════════

Aviation engineering has achieved an extraordinary level of mechanical reliability
over the course of its history. According to data from the International Civil Aviation
Organization (ICAO), the accident rate per million flights in commercial aviation has
decreased approximately fortyfold since the early 1970s \citep{icao2022}. Yet behind
this remarkable achievement in mechanical safety, a structural gap persists that has not
been fully resolved: the inability to explain, at a diagnostic level, the human factors
underlying aviation accidents.

A review of the aviation safety literature reveals that existing approaches have coalesced
around two fundamental paradigms. The reactive paradigm seeks to prevent the recurrence of
similar events by conducting systematic analyses following accidents and incidents \citep{gerede2018}.
The proactive paradigm aims to identify risks before they materialize, working to elevate
operational safety through tools such as Safety Management Systems (SMS) and the Human
Factors Analysis and Classification System (HFACS) \citep{shappell2000}. While the contributions
of these two paradigms to aviation safety are undeniable, both are structurally incapable of
objectively recording the pilot's cognitive and physiological state at the moment of an accident.

The traditional accident investigation tools---the FDR and CVR---capture the aircraft's
mechanical responses and cockpit conversations. Yet what the pilot was experiencing at those
critical moments, the level of oxygen saturation, or how heart rate variability was evolving
under acute stress, remains largely in the dark \citep{dehais2020}. This situation creates
fertile ground for the proliferation of superficial causal attributions such as ``pilot error''
in accident reports---attributions that simultaneously conceal genuine root causes and render
systemic failures invisible \citep{reason2000}.

Recent advances in wearable sensor technologies carry considerable potential for closing this
diagnostic gap. Smart pilot watches and biometric sensors can measure parameters such as HRV,
SpO\textsubscript{2}, and GSR in real time; when fused with traditional flight telemetry,
these data offer the possibility of reconstructing the physiological state at the moment of
an accident with mathematical precision \citep{cantekin2025a}. However, the obstacles to
realizing this technological potential are less technical than ethical and legal. A pilot's
biometric data---heart rate, stress level, sleep quality---are classified as special category
personal data under the European Union's General Data Protection Regulation (GDPR), and
questions of ownership and access rights over such data remain unresolved critical issues
\citep{erdinc2020}.

This study proposes the Complementary Diagnostic Safety Approach (CDSA)---a framework that
addresses these ethical and technical challenges simultaneously---and presents in detail the
cryptographic infrastructure underpinning this approach. CDSA is positioned as a third safety
layer that goes beyond reactive and proactive paradigms. Its core innovation lies in a two-phase
data processing architecture based on AES-256-GCM encryption and Shamir Secret Sharing protocol.
Through this architecture, individual biometric data are protected by cryptographic guarantees
while anonymized macro data contribute to proactive risk management.

The original contributions of this study can be summarized across four dimensions:
(1) the reconceptualization of aviation safety within a three-layer paradigm framework
encompassing reactive, proactive, and diagnostic approaches;
(2) the design of a cryptographic architecture integrating AES-256-GCM and Shamir Secret
Sharing with Just Culture principles;
(3) the adaptation of differential privacy mechanisms to aviation safety management; and
(4) the operationalization of the CDSA ecosystem through publicly accessible prototype
applications on the \href{https://cdsa.app}{cdsa.app} platform.
The remainder of the paper is organized as follows: Section~\ref{sec:related} reviews the
relevant literature; Section~\ref{sec:framework} presents the CDSA framework and cryptographic
architecture; Section~\ref{sec:prototype} reports prototype implementation and validation findings;
Section~\ref{sec:discussion} discusses ethical and regulatory dimensions; and
Section~\ref{sec:conclusion} evaluates the study's contributions and future research directions.

% ══════════════════════════════════════════════════════════════════════════
\section{Related Work}
\label{sec:related}
% ══════════════════════════════════════════════════════════════════════════

The following five axes of literature are reviewed: the evolution of safety paradigms in
aviation, wearable biometric sensor applications, cryptographic privacy architectures, data
privacy and Just Culture principles, and the identification of the gap that this study addresses.

\subsection{The Evolution of Paradigms in Aviation Safety}

A field's understanding of safety is directly linked to how the accidents within that field
are interpreted. A review of the aviation safety literature reveals that the discipline's
history of more than a century has been built on a reactive foundation. Early accident
investigations relied on linear causality models in which events were attributed to a single
cause---typically a mechanical failure or human error. \citet{reason2000}'s Swiss Cheese model
fundamentally transformed this understanding, demonstrating that accidents arise from the
simultaneous alignment of gaps across multiple defensive layers. This model established systemic
thinking in aviation and formed the theoretical basis for proactive safety management.

The operational counterpart of proactive safety approaches is ICAO's Safety Management System
(SMS) framework \citep{icao2020}. SMS seeks to systematically integrate hazard identification,
risk assessment, and safety assurance cycles into institutional operations. The HFACS model
developed by \citet{shappell2000} classifies human factor-related accidents within a four-level
hierarchical structure, providing accident investigators with a powerful analytical tool. Despite
these advances, \citet{gerede2018} emphasizes that existing proactive tools remain structurally
limited in their diagnostic capacity precisely because they cannot access the physiological data
layer at the moment of an accident.

\subsection{Wearable Biometric Sensors and Aviation Applications}

The potential of wearable technologies in the aviation context has attracted growing interest
over the past decade. Heart rate variability (HRV) serves as a direct indicator of autonomic
nervous system activity, objectively reflecting the pilot's cognitive workload and stress level
\citep{taskforce1996}. \citet{dehais2020}, in a study conducted from a neuroergonomics
perspective, demonstrated that HRV reduction reliably predicts prefrontal cortex overload and
consequent loss of situational awareness. Peripheral oxygen saturation (SpO\textsubscript{2})
functions as a critical biomarker for the early detection of covert hypoxia risk, particularly
under high-altitude conditions. \citet{denison1972} demonstrated that the pilot's reasoning
capacity begins to deteriorate well before SpO\textsubscript{2} drops below 91\,\%.

On the integration of biometric data with flight safety, \citet{ratti2021} showed that sensor
fusion algorithms, when combined with FDR telemetry, can reconstruct the causal chain at the
moment of an accident with mathematical accuracy. \citet{cantekin2025a} comprehensively examined
the diagnostic potential of pilot watches evolving from the mechanical era to AI-assisted systems
and laid out the first systematic framework of the CDSA concept.

\subsection{Cryptographic Privacy Architectures}

The most critical legal and ethical dimension of biometric data processing is the provision
of data privacy guarantees at the design level. \citet{cavoukian2009}'s Privacy by Design
principle advocates that personal data protection mechanisms be conceived as an integral part
of the architecture rather than added to the system after the fact. This principle has been
elevated to the status of a legal obligation under Article~25 of the GDPR \citep{gdpr2016}.

In the field of threshold cryptography, \citet{shamir1979} developed a protocol enabling a
secret value to be split into $n$ shares such that any $k \geq 2$ shares can reconstruct it
via Lagrange interpolation. Differential privacy, formally defined by \citet{dwork2014}, renders
individual re-identification statistically impossible by adding controlled random noise to a dataset.
The application of cryptographic access control and threshold schemes to aviation data management
finds very limited coverage in the literature; the NTSB's \citeyearpar{ntsb2023} projection report
explicitly articulates the need for a regulatory framework in this area.

The relationship between regulatory complexity and cryptographic data protection architecture is
not unique to the aviation sector. In healthcare, HIPAA and GDPR mandate that patient data be
protected through encryption, access control, and audit trail mechanisms; compliance must be
grounded in direct technical safeguards rather than institutional trust alone \citep{conduah2025}.
In the financial sector, SOX, MiFID~II, and Basel~III frameworks require comprehensive data
governance infrastructures, and the compliance burden simultaneously creates structural barriers
to market entry for new technological actors \citep{mcnulty2023}. In defense and intelligence,
ITAR and FedRAMP~High certifications demand compliance processes spanning more than a decade.
A common pattern emerges: regulatory cryptographic obligations strengthen the market position
of compliant platforms; privacy-enhancing technologies are increasingly regarded not merely
as compliance tools but as prerequisites for value creation in the data economy \citep{nitrd2022}.
CDSA represents a sector-specific application of this cross-industry data protection paradigm.

\subsection{Data Privacy, Just Culture, and Pilot Rights}

The point at which biometric monitoring systems intersect with pilot rights constitutes the most
contested axis in the literature. The Just Culture principle conceptualized by \citet{dekker2007}
holds that errors should be reported not to be punished but to identify vulnerabilities in the system.
The Air Line Pilots Association \citeyearpar{alpa2022} documented in its technical white paper that
the absence of this guarantee would fundamentally undermine pilot trust and produce a counterproductive
effect on safety culture. \citet{ustaomer2020} demonstrated that the institutional trust environment
is the single most critical factor determining safety reporting behavior. \citet{erdinc2020}
evaluated the legal requirements for processing biometric data under the special category personal
data framework through the lenses of proportionality and informed consent.

\subsection{The Gap in the Literature and the Position of This Study}

Two fundamental observations emerge from the body of literature summarized above. First, studies
addressing the integration of wearable biometric sensors into aviation diagnostic processes discuss
the technical opportunity framework but fall short of proposing an end-to-end cryptographic
architecture to secure this integration. Second, the data privacy and Just Culture literature
evaluates the ethical framework broadly but cannot systematically demonstrate how these principles
translate into engineering design. This study aims to close these gaps through the integrated
application of AES-256-GCM real-time encryption, Shamir 2-of-3 threshold cryptography, and
Laplace differential privacy.

% ══════════════════════════════════════════════════════════════════════════
\section{CDSA Framework and Cryptographic Architecture}
\label{sec:framework}
% ══════════════════════════════════════════════════════════════════════════

This section presents the conceptual foundation of CDSA and its two-phase cryptographic
architecture in detail. Phase~1 addresses individual diagnostic access control;
Phase~2 covers anonymized macro-level data analysis for proactive risk management.

\subsection{Conceptual Framework: The Third Safety Paradigm}

A system's diagnostic capacity is determined not only by the effectiveness of its preventive
mechanisms but also by the extent to which it can illuminate the causal chain underlying events.
Viewed from this perspective, the current reactive and proactive paradigms in aviation safety
share a common blind spot: both are structurally incapable of objectively recording the pilot's
physiological and cognitive state at the moment of an accident.

The CDSA proposed in this study is positioned as a third safety layer targeting precisely this
blind spot. The core thesis of CDSA can be stated as follows: when biometric data obtained from
wearable sensors are fused with traditional flight telemetry at millisecond precision, accident
investigations can move beyond superficial attributions such as ``pilot error'' and physiological
root causes---covert hypoxia, acute stress, or startle effect---can be established with
mathematical certainty. CDSA seeks to resolve the tension between individual privacy and diagnostic
access not through legal regulation but through the architecture of cryptography itself.

\subsection{Phase~1: Cryptographic Diagnostic Data Architecture}

Phase~1 comprises a five-layer hierarchical architecture designed to make individual pilot
biometric data accessible only following an accident or serious incident report and subject
to simultaneous multi-party authorization.

\subsubsection{System Overview}

Phase~1 of CDSA is built on a five-layer hierarchical architecture: (1)~real-time biometric
measurement and encryption, (2)~splitting of the master key into shares via threshold
cryptography, (3)~distribution of key shares to independent institutional parties,
(4)~unlocking of the locked box following an accident report and with threshold authorization,
and (5)~decryption and sensor fusion with FDR/CVR data. Figure~\ref{fig:architecture}
summarizes the architecture and data flow across these five layers.

\begin{figure}[H]
  \centering
  \includegraphics[width=\textwidth]{image1.png}
  \caption{CDSA Two-Phase Cryptographic Architecture. Phase~1 (left): AES-256-GCM
  encryption, Shamir 2/3 threshold key distribution, and locked-box access control.
  Phase~2 (right): Differential privacy anonymization and SMS integration. Both phases
  share the same hardware infrastructure while maintaining cryptographic separation of access.}
  \label{fig:architecture}
\end{figure}

\subsubsection{Data Acquisition and Real-Time Encryption}

The most effective form of data security is considered to be ensuring that sensitive information
is never stored in plaintext. Accordingly, the design principle adopted is that HRV,
SpO\textsubscript{2}, and GSR values continuously acquired from the smart pilot watch at
millisecond precision are never written to device memory in raw form. Each measurement is
encrypted in real time at the hardware level using the AES-256-GCM algorithm (NIST SP 800-38D).
The 96-bit nonce value used in encryption is generated independently for each record by a
cryptographically secure pseudorandom number generator (CSPRNG), ensuring that identical data
never produce the same ciphertext. Each record additionally carries an Additional Authentication
Data (AAD) field composed of the flight identifier, timestamp, and box identifier; this field
is cryptographically bound to the GCM authentication tag.

\subsubsection{Shamir Secret Sharing and 2-of-3 Threshold Cryptography}

To eliminate the security vulnerability created by a single center's absolute control over the
master key, the AES-256 master key ($k_0$) is split into $n=3$ shares using the Shamir Secret
Sharing protocol \citep{shamir1979}. The mathematical foundation of the system is expressed as:
%
\begin{equation}
  f(x) = k_0 + a_1 x \pmod{p}, \quad p = 2^{257} - 93
  \label{eq:shamir_poly}
\end{equation}
\begin{equation}
  S_i = \bigl(x_i,\, f(x_i)\bigr) \;\longrightarrow\; \text{Share } i
  \text{ transmitted to Party } i
  \label{eq:shamir_shares}
\end{equation}
%
Here $p$ is a 257-bit prime compliant with NIST SP 800-57 standards. The threshold parameter
$k=2$ means that any two of the three shares suffice to reconstruct the master key exactly
via Lagrange interpolation. A single share held by one party is mathematically meaningless
and cannot yield the original key even with unlimited computational power.
Table~\ref{tab:crypto_params} summarizes the cryptographic parameters of the CDSA architecture.

\begin{table}[H]
  \centering
  \caption{CDSA Cryptographic Architecture Parameter Summary.}
  \label{tab:crypto_params}
  \begin{tabularx}{\textwidth}{lX}
    \toprule
    \textbf{Parameter} & \textbf{Value / Description} \\
    \midrule
    Symmetric encryption      & AES-256-GCM (NIST SP 800-38D) \\
    Nonce length              & 96-bit (CSPRNG, independent per record) \\
    Secret sharing protocol   & Shamir Secret Sharing \citep{shamir1979} \\
    Finite field              & $p = 2^{257} - 93$ (257-bit prime, NIST SP 800-57) \\
    Threshold parameters      & $k=2$, $n=3$ (2-of-3 threshold scheme) \\
    AAD field                 & Flight ID + Timestamp + Box ID \\
    Integrity control         & GCM authentication tag \\
    \bottomrule
  \end{tabularx}
\end{table}

Within the parameter framework summarized in Table~\ref{tab:crypto_params},
Figure~\ref{fig:shamir} illustrates the protocol flow through which Shamir Secret Sharing
splits the master key into three independent shares and reconstructs it when the threshold
number of shares is combined.

\begin{figure}[H]
  \centering
  \includegraphics[width=\textwidth]{image2.png}
  \caption{Shamir Secret Sharing: 2-of-3 Threshold Cryptography Protocol. The AES-256
  master key $k_0$ is split into $n=3$ shares using a polynomial over a finite field
  ($p=2^{257}-93$). Any 2 shares allow reconstruction via Lagrange interpolation; a
  single share alone reveals zero information.}
  \label{fig:shamir}
\end{figure}

\subsubsection{Independent Key Holders and Institutional Balance}

The effectiveness of cryptographic systems is measured not only by their mathematical strength
but also by the institutional independence of their key holders. Accordingly, the three key
shares are distributed to three independent institutional parties whose interests are structurally opposed:

\begin{itemize}[leftmargin=1.5cm]
  \item \textbf{Share A --- Accident Investigation Authority (NTSB/JTSB):} Submits the key
    share upon completion of accident reporting. Its legal status, independent of commercial
    pressures, constitutes the primary safeguard.
  \item \textbf{Share B --- Pilot Union (IFALPA/ALPA):} Represents pilot consent in its
    capacity as the guardian of labor rights. Retains veto power over the use of data in
    disciplinary proceedings.
  \item \textbf{Share C --- Judicial Authority:} Activated solely upon court order and upon
    presentation of legal justification.
\end{itemize}

This tripartite structure renders it mathematically impossible for any single party to unlock
the system alone, while simultaneously aligning with Just Culture principles \citep{dekker2007}.

\subsubsection{Locked Box Access Control and Decryption Protocol}

During routine flight operations, encrypted biometric data accumulate in hardware units referred
to as the locked box. After the flight concludes and the box is sealed, the master key is
cryptographically deleted from device memory. At this stage, even if the airline or aircraft
manufacturer physically possesses the encrypted data, they cannot access its contents. The
decryption mechanism can only be triggered by the simultaneous fulfillment of two conditions:
the filing of an accident or serious incident report, and the simultaneous submission of key
shares by at least two parties. The decrypted HRV, SpO\textsubscript{2}, and GSR time series
are synchronized with FDR telemetry at millisecond precision and presented to the accident
investigation board.

\subsection{Phase~2: Macro Data Analysis with Differential Privacy}

Phase~2 shares the same biometric hardware infrastructure as Phase~1 but is fully decoupled from
it through cryptographic access control. In this phase, anonymized macro data containing no
individual identity are processed for proactive risk management.

\subsubsection{Anonymization Architecture}

While Phase~1 of CDSA serves a diagnostic function, Phase~2 leverages the same biometric data
infrastructure for proactive risk management purposes. Data stripped of individual pilot identity
pass through two complementary privacy mechanisms. The first mechanism is based on the
$k$-anonymity ($k=5$) principle \citep{sweeney2002}. The second mechanism is Laplace
differential privacy ($\varepsilon=1.0$) \citep{dwork2014}:
%
\begin{equation}
  \mathcal{M}(x) = x + \mathrm{Lap}\!\left(\frac{\Delta f}{\varepsilon}\right),
  \quad \Delta f = \text{sensitivity},\quad \varepsilon = \text{privacy budget}
  \label{eq:dp}
\end{equation}
%
The sensitivity values applied are $\Delta f = 5.0$\,ms for HRV, $\Delta f = 1.0\%$ for
SpO\textsubscript{2}, and $\Delta f = 2.0\,\mu$S for GSR.

\subsubsection{SMS Integration and Proactive Risk Outputs}

Anonymized macro data are integrated with SMS to produce two types of operational output:
(1)~evidence-based justification for shift optimization in route and shift combinations where
high SRS values are concentrated; and (2)~support for ICAO Annex~6 FDP limitations through the
finding that night shifts produce statistically significantly higher risk profiles across all routes.
In neither output does individual pilot data enter the processing pipeline.

\subsection{Comparative Summary of the Architecture}

Figure~\ref{fig:comparison} below presents a visual comparative summary of the two phases across
data type, encryption mechanism, access condition, authorized users, primary goal, and GDPR basis.

\begin{figure}[H]
  \centering
  \includegraphics[width=\textwidth]{image3.png}
  \caption{CDSA Phase~1 vs Phase~2: Comparative Data Flow. Key distinctions across data type,
  privacy mechanism, access trigger, authorized users, primary goal, GDPR basis, and operational
  output. Both phases share the same biometric hardware infrastructure.}
  \label{fig:comparison}
\end{figure}

% ══════════════════════════════════════════════════════════════════════════
\section{Prototype Implementation and Validation}
\label{sec:prototype}
% ══════════════════════════════════════════════════════════════════════════

This section validates the functionality of the CDSA architecture at two complementary levels:
a Python-based cryptographic prototype implementation and publicly accessible user interfaces
deployed on the \href{https://cdsa.app}{cdsa.app} platform.

\subsection{Validation Strategy and Prototype Scope}

The scientific value of a theoretical framework depends in large part on the extent to which
its practical feasibility can be concretely demonstrated. The former establishes mathematical
correctness; the latter demonstrates operational feasibility.

A synthetic covert hypoxia scenario formed the basis of the prototype development process:
SpO\textsubscript{2}\,=\,91\%, HRV\,=\,33.2\,ms, GSR\,=\,13.96\,$\mu$S, SRS\,=\,0.457
(above the critical threshold of 0.40). These parameters are consistent with the
literature-calibrated synthetic dataset (seed\,=\,42).

\subsection{Cryptographic Prototype --- Python Implementation}

The cryptographic prototype was developed and validated at two levels: the implementation
environment and components are described in Section~\ref{sec:421}, while the Phase~1 validation
findings are reported in Section~\ref{sec:422}.

\subsubsection{Implementation Environment and Components}
\label{sec:421}

\sloppy The cryptographic prototype was developed in Python~3.10 using the \texttt{cryptography} library for encryption operations. The prototype consists of three modules: (1)~data generation
and real-time encryption, (2)~key management via Shamir Secret Sharing, and (3)~decryption and
integrity verification via Lagrange interpolation. The complete code is presented in Appendix~C
and shared publicly in accordance with the reproducibility principle.

\subsubsection{Phase~1 Validation Findings}
\label{sec:422}

In the encryption module, six biometric measurement records from the hypoxia scenario defined
above were encrypted in real time using AES-256-GCM. Independent 96-bit nonce values generated
for each record provided mathematical protection against timing attacks. Successful binding of
AAD fields to GCM authentication tags was confirmed.

In the Shamir Secret Sharing module, the master key $k_0$ was split into $n=3$ shares over the
finite field $p=2^{257}-93$, with each share transmitted to the corresponding institutional party.
It was verified that using only two of the three shares---NTSB and IFALPA, without the judicial
authority's share---the master key was fully reconstructed via Lagrange interpolation. The key
integrity verification test returned a positive result; single-party access attempts were
automatically rejected.

The decryption output reconstructed the covert hypoxia scenario with mathematical precision:
SpO\textsubscript{2} had dropped to 91\%, HRV had declined to 33.2\,ms, and GSR had risen to
13.96\,$\mu$S at the moment of critical event onset---findings that establish the existence of
a physiological root cause chain beyond the attribution of ``pilot error.''
Figure~\ref{fig:pilotauth} presents the PilotSense AUTH interface alongside the decrypted
biometric evidence timeline for the TC-XXX scenario.

\begin{figure}[H]
  \centering
  \includegraphics[width=\textwidth]{image4.png}
  \caption{PilotSense AUTH Interface: Phase~1 Diagnostic Access Control. Left: encrypted
  locked-box with NTSB and IFALPA shares submitted (2/3 threshold reached). Right: decrypted
  biometric evidence timeline revealing covert hypoxia onset 28 minutes before the reported
  event (synthetic scenario, TC-XXX).}
  \label{fig:pilotauth}
\end{figure}

\subsection{Prototype Applications --- Overview}

The CDSA ecosystem comprises five prototype applications, each targeting a distinct operational
layer. Table~\ref{tab:apps} summarizes the key metrics obtained for each application.

\begin{table}[H]
  \centering
  \caption{Summary of CDSA Prototype Applications.}
  \label{tab:apps}
  \begin{tabularx}{\textwidth}{llX}
    \toprule
    \textbf{Application} & \textbf{Layer} & \textbf{Key Metric} \\
    \midrule
    PilotSense AUTH & Diagnostic --- Phase~1 & 2/3 multi-party signature, AES-256-GCM, GDPR Art.~25 \\
    PilotSense OPS  & Proactive --- Phase~2 & 247 pilots, SMS integration, anonymous fleet monitoring \\
    PilotO2         & Proactive --- Individual & 95\% model accuracy, 92\% hypoxia detection, 127\,ms \\
    PilotGuard      & Proactive --- Individual & 87\% pre-flight risk, 93\% stress detection \\
    PilotReflect    & Post-Flight (EFB)       & 8-min debriefing, ISO 9241-11, Privacy by Design \\
    \bottomrule
  \end{tabularx}
\end{table}
\fussy

All applications are publicly accessible on the \href{https://cdsa.app}{cdsa.app} platform.
It must be emphasized that all five applications remain at the prototype stage and FAA/EASA
certification processes have not yet been completed.

\subsubsection{PilotSense AUTH --- Diagnostic Layer Interface}
PilotSense AUTH operationalizes Phase~1 of CDSA at the user interface level. In the TC-XXX
simulation scenario, a 14.7\,MB encrypted biometric data file was successfully decrypted through
the simultaneous submission of NTSB and IFALPA shares, demonstrating the practical functionality
of the 2/3 threshold scheme. The GDPR Article~25 compliance layer remained active throughout.

\subsubsection{PilotSense OPS --- Proactive Fleet Monitoring Interface}
PilotSense OPS represents Phase~2 of CDSA, targeting airline SMS managers and operations centers.
In the simulation environment, anonymized biometric data from 247 active pilots were successfully
processed; eight warnings and two critical conditions were automatically detected. The SRS values
of critical records exceeded the threshold of 0.40.

\subsubsection{PilotO2 --- Covert Hypoxia Early Warning System}
The PilotO2 application is built on a three-layer LSTM neural network architecture consisting
of 128, 64, and 32 processing units. The system achieved 95\% model accuracy, 92\% hypoxia
detection rate, and 127\,ms response time. Validation under real operational conditions
constitutes a priority for future research.

\subsubsection{PilotGuard --- Pre-Flight Readiness Assessment}
PilotGuard processes HRV, instantaneous heart rate, stress index, and sleep quality data in
real time to generate an integrated pilot flight readiness score. Simulation tests yielded
87\% pre-flight risk prediction and 93\% stress detection accuracy, with a user experience
rating of 4.6/5.0.

\subsubsection{PilotReflect --- EFB Integrated Biometric Debrief System}
PilotReflect delivers post-flight personal biometric assessment within the Type~B EFB ecosystem.
In line with \citeauthor{cavoukian2009}'s \citeyearpar{cavoukian2009} Privacy by Design
principles, all biometric data are stored exclusively on the pilot's personal device encrypted
with AES-256. An average debriefing duration of eight minutes and an FAA AC 120-76E compliant
interface design were achieved in accordance with ISO 9241-11 usability criteria.

\subsection{Integrated Assessment of Prototype Findings}

Viewed collectively, the prototype validation findings demonstrate that the CDSA ecosystem can
be successfully translated from a theoretical framework into a functional platform. Several
limitations must be carefully underscored. All five applications remain at the prototype stage;
FAA/EASA certification processes have not yet been completed. The cryptographic validation was
conducted on a synthetic dataset, and sensor behavior under real operational conditions---G-forces,
airframe vibration---has not yet been tested.

% ══════════════════════════════════════════════════════════════════════════
\section{Discussion}
\label{sec:discussion}
% ══════════════════════════════════════════════════════════════════════════

This section examines the ethical coherence of the cryptographic architecture, the Panopticon
paradox and its design solution, compliance with GDPR and the international regulatory
framework, the limitations of synthetic data, comparative positioning in the literature,
and practical feasibility.

\subsection{Ethical Coherence of the Cryptographic Architecture}

For a technological solution to be sustainable, it must satisfy not only engineering criteria
but also the value systems of the communities it affects. The most distinctive contribution of
the CDSA architecture is considered to lie not in its technical innovation but in its ethical
coherence. Existing aviation data systems attempt to manage the tension between privacy and
diagnostic access through legal regulation---an approach whose guarantees shift whenever
legislation changes or interpretations diverge. CDSA transfers this tension to the mathematical
foundations of cryptography, thereby eliminating this fragility.

The alignment of the 2-of-3 threshold scheme with Just Culture philosophy is particularly
noteworthy in this context. In the CDSA architecture, the pilot union's possession of a key
share prevents biometric data from being used in disciplinary proceedings not through a legal
guarantee but through a mathematical impossibility. As \citet{alpa2022} emphasizes, pilot
organizations systematically oppose biometric monitoring systems precisely because of their
distrust of existing legal safeguards---making the mathematical nature of the guarantee, rather
than its legal form, the critical differentiator.

\subsection{The Panopticon Paradox and the Design Solution}

The most insidious risk created by continuous biometric monitoring is not the direct harm it
may cause to the pilot but the behavioral change it may induce. \citeauthor{foucault1992}'s
\citeyearpar{foucault1992} Panopticon metaphor illustrates this mechanism compellingly: an
individual who knows they are under constant surveillance begins to behave as though being
watched even in the absence of observation. Transposed to the aviation context, this effect
carries the potential to amplify performance anxiety, elevate allostatic load, and paradoxically
threaten flight safety \citep{zuboff2019}.

CDSA's response to this paradox is direct and structural. During routine flight operations,
airline management mathematically cannot access individual biometric data. The pilot knows that
the system is not monitoring them at the individual level; they are only aware that their
anonymized macro data contribute to fleet-wide risk patterns. This design choice is considered
to eliminate the Panopticon effect. It must be acknowledged, however, that this guarantee holds
only as long as the system functions correctly; a security breach or data leak could
fundamentally destabilize this balance.

\subsection{Compliance with GDPR and the International Regulatory Framework}

The legal positioning of the CDSA architecture simultaneously engages multiple regulatory layers.
The individual biometric data processed in Phase~1 are classified as special category personal
data under GDPR Article~9 \citep{gdpr2016}. CDSA addresses this requirement by integrating
pilot consent into key management through Share~B---an approach that functions as both a technical
and a legal safeguard. The anonymized macro data in Phase~2 fall outside the scope of GDPR's
definition of personal data following the application of differential privacy mechanisms.
At the international level, a Biometric Flight Data Declaration is presented as this study's
concrete policy proposal for updating the ICAO Annex~13 framework \citep{icao2020}.

\subsection{Limitations of Synthetic Data and Ecological Validity}

Addressing the methodological limitations of this study openly and honestly is a fundamental
requirement of academic integrity. The biometric data used to validate the cryptographic
architecture were derived from a literature-calibrated synthetic dataset ($n=1{,}179$,
seed\,=\,42). This dataset has established its statistical validity through Shapiro-Wilk
subgroup normality testing (13/16 subgroups, $p>0.05$), Spearman correlation ($\rho=-0.648$,
$p<0.001$), and Mann-Whitney~U ($p<0.001$) tests; however, its ecological validity under real
cockpit conditions has not yet been demonstrated. Figure~\ref{fig:heatmap} presents the
HRV--GSR correlation structure and the route-by-shift SRS risk heatmap derived from this
synthetic dataset.

\begin{figure}[H]
  \centering
  \includegraphics[width=\textwidth]{image5.png}
  \caption{HRV\,\textasciitilde\,GSR Correlation and SRS Risk Heatmap (Synthetic Dataset,
  $n=1{,}179$, seed\,=\,42). Left: Spearman $\rho=-0.648$, $p<0.001$, consistent with
  \citet{dehais2020}. Right: Route\,$\times$\,Shift SRS heatmap. IST-SYD Night
  (SRS\,=\,0.440) approaches the critical threshold of 0.40.
  Source: Literature-calibrated synthetic data (Appendix~B).}
  \label{fig:heatmap}
\end{figure}

Three ecological validity concerns are identified. First, the extent to which high G-forces and
airframe vibration introduce noise into biometric sensor measurements is unknown. Second, variance
in individual pilot baseline values is not represented in the synthetic data, rendering it
impossible to personalize the SRS critical threshold on a per-pilot basis. Third, the impact of
sensor calibration drift on system performance during extended operational use has not yet been
modeled. Addressing these limitations constitutes the primary objective of the three-phase
validation plan presented in Section~\ref{sec:63}.

\subsection{Comparative Positioning: CDSA in the Literature}

The CDSA architecture sits at the intersection of three related but independently evolving
streams of research: physiological monitoring studies in aviation safety literature, threshold
scheme applications in applied cryptography, and differential privacy systems in data privacy
literature. \citet{dehais2020} comprehensively documented the diagnostic potential of biometric
monitoring yet did not address cryptographic privacy guarantees. \citet{dwork2014} established
the theoretical foundations of differential privacy on solid ground but did not discuss
operational applications in high-risk domains such as aviation. The NTSB's \citeyearpar{ntsb2023}
projection report proposed the integration of wearable device data into accident investigations
yet was unable to offer any technical framework for how this integration should be supported by
cryptographic guarantees. CDSA aims to close these gaps within a single unified architecture.

Beyond the literature gaps identified above, the cryptographic guarantee approach proposed in
this study represents a sector-specific application of data protection paradigms long established
in other highly regulated industries. In healthcare, HIPAA and GDPR mandate technically grounded
data protection through encryption, access control, and audit trail mechanisms \citep{conduah2025}.
In finance, SOX, MiFID~II, and Basel~III frameworks impose comprehensive data governance
obligations that simultaneously create structural entry barriers for new market entrants
\citep{mcnulty2023}. Across these sectors, a consistent pattern emerges: regulatory cryptographic
obligations strengthen the market position of compliant platforms, and privacy-enhancing
technologies are increasingly regarded not merely as compliance tools but as prerequisites for
value creation in the data economy \citep{nitrd2022}. CDSA, positioned as the aviation-specific
instantiation of this cross-industry paradigm, inherits both the methodological legitimacy and
the strategic positioning that this broader pattern affords.

\subsection{Practical Feasibility and Operational Barriers}

The first barrier is certification cost. The DO-178C Level~B compliance process demands
significant financial and temporal resources. Smaller and medium-sized airlines may struggle to
bear this burden independently, potentially delaying sector-wide adoption. Shared certification
pools and industry consortia represent viable pathways for overcoming this challenge.

The second barrier is institutional cultural resistance. Pilot organizations harbor deep and
legitimate reservations about biometric monitoring \citep{alpa2022}. Although the CDSA 2-of-3
threshold scheme renders disciplinary use of data mathematically impossible, translating
technical guarantees into institutional trust requires sustained transparency policies and
long-term stakeholder dialogue.

The third barrier is the requirement for multi-party coordination. The legal framework governing
the simultaneous participation of the accident investigation authority, pilot union, and judicial
authority during international flights has not yet been clarified. Incorporation of a ``Personal
Biometric Flight Data'' standard provision into ICAO Annex~13 represents the most operationally
feasible solution to this gap.

The fourth barrier is the absence of real-world data. All cryptographic validation work in this
study was conducted on synthetic data. The effects of real cockpit conditions---G-forces,
airframe vibration, and inter-pilot baseline variability---on system performance remain unknown.
This uncertainty limits the confidence of regulators and pilot organizations in the system,
making the three-phase validation roadmap presented in Section~\ref{sec:63} a strategic imperative.

% ══════════════════════════════════════════════════════════════════════════
\section{Conclusion}
\label{sec:conclusion}
% ══════════════════════════════════════════════════════════════════════════

This section summarizes the original contributions of the study across four dimensions, states
the methodological limitations explicitly, and identifies three priority directions for future research.

\subsection{Assessment of Original Contributions}

The true maturity of a safety system is measured not only by its capacity to prevent events but
also by its ability to fully comprehend the human reality underlying those events. The CDSA
framework proposed in this study has been positioned as a third layer complementing the reactive
and proactive paradigms of aviation safety, and an integrated cryptographic architecture has been
developed capable of illuminating---with mathematical precision---the root causes of human factors
in accident investigations.

The four original contributions of this study can be summarized as follows. The first contribution
is conceptual: the systematic introduction of the diagnostic paradigm into the literature alongside
the reactive and proactive paradigms, carrying the potential to fundamentally transform how human
factors are approached in accident investigations. The second contribution is technical: the
integration of AES-256-GCM and Shamir 2-of-3 threshold cryptography with Privacy by Design
principles, grounding biometric data privacy in a mathematical foundation independent of legal
guarantees. The third contribution is methodological: the integration of differential privacy
mechanisms ($\varepsilon=1.0$, $k$-anonymity $k=5$) with Safety Management Systems, resolving
the tension between individual privacy and collective risk management at the statistical level.
The fourth contribution is applied: the operationalization of the CDSA ecosystem through publicly
accessible prototype applications on the \href{https://cdsa.app}{cdsa.app} platform, demonstrating
the framework's feasibility beyond abstract proposal. Furthermore, the cryptographic guarantee
approach proposed in this study represents a sector-specific application of data protection
paradigms long established in healthcare (HIPAA), finance (SOX/MiFID~II), and defense
(ITAR/FedRAMP) --- positioning CDSA as the aviation instantiation of a cross-industry regulatory
data protection pattern \citep{conduah2025,mcnulty2023,nitrd2022}.

\subsection{Limitations}

All cryptographic validation findings in this study were conducted on a literature-calibrated
synthetic dataset ($n=1{,}179$, seed\,=\,42). Completing DO-178C Level~B or higher compliance
before CDSA components can be operationally deployed does not appear feasible within the current
regulatory framework. The legal framework governing the simultaneous participation of the accident
investigation authority, pilot union, and judicial authority in international flights has not yet
been clarified. Beyond these structural limitations, it must be explicitly acknowledged that the
ecological validity of the synthetic dataset findings under real cockpit conditions---including
the effects of G-forces, airframe vibration, and inter-pilot baseline variability---remains
unestablished.

\subsection{Future Research Roadmap}
\label{sec:63}

Three priority research directions emerge from a joint assessment of this study's findings and
limitations. Real operational data validation is the first priority: the collection of real
operational biometric data through voluntary pilot participation and in collaboration with
IFALPA/ALPA, and the application of the same regression framework, are planned. This validation
process is envisioned in a three-phase structure: (1)~establishment of ethical approval and data
access protocols, (2)~Anderson-Darling distribution fit testing and comparison of SRS patterns
with real data, and (3)~updating of the calibration matrix to transform CDSA into an operational
tool. Initiating the FAA/EASA certification process constitutes the second priority. ICAO
Annex~13 integration---through the advancement of the Biometric Flight Data Declaration into the
policy process---constitutes the third priority.

\subsection{Closing: Toward a Sustainable Balance Between Privacy and Diagnostic Power}

The future of aviation safety most likely lies in a better understanding of human beings rather
than machines. While aircraft have approached a near-perfect level of mechanical reliability,
the human factor remains the most complex and least understood component of the system. The
wearable technology ecosystem carries the potential to transform this equation---yet this
potential must be realized without compromising pilot rights, undermining safety culture, or
eroding institutional trust.

The CDSA framework proposed in this study seeks to resolve the tension between privacy and
diagnostic power through the design of technology itself. The system's success ultimately depends
on the simultaneous fulfillment of three conditions: validation of the cryptographic architecture
against real operational data, updating of the regulatory framework to accommodate this innovation,
and the establishment of long-term dialogue with pilot organizations. To the extent that these
three conditions can be met concurrently, CDSA will be capable of substantiating its claim to be
the third paradigm of aviation safety.

% ── References ────────────────────────────────────────────────────────────
\bibliographystyle{apalike}
\bibliography{references}

\end{document}
